\section{Recovering a local density from five auxiliary periods}
\label{sec:five-period-localization}
The local projection in the nonlinear step uses five small periodic
boxes. Their side lengths differ by two. A fourth finite difference
cancels all the contributions from folded boundary faces and retains
the interior contribution. We give the construction, since both the
exact identity and its derivative cost matter for the parameter
estimates.

\subsection{Folded fields and boundary types}
For an integer $w\ge1$, fold the cyclic set $\mathbb Z/(2w)\mathbb Z$
onto $\{0,\ldots,w\}$ by
\[
       \pi_w(q)=\min(\bar q,2w-\bar q),\qquad 0\le\bar q<2w.
\]
Its two endpoints have one preimage each and its other points have
two. Apply this map in each of the four coordinates. Adjacent sites
map to adjacent sites, with direction either preserved or reversed.
For a gauge field $U$ on a larger torus and an anchor $o$, the
folded field $\Phi_{w,o}U$ on the side-$2w$ torus reads the link
of $U$ between $o+\pi_w(q)$ and $o+\pi_w(q+e_\mu)$; a reversed
link is read as its inverse. Here and below a lattice point is
reduced modulo the large period when used to index $U$.

Fix $r\ge1$ and $w\ge2r$. In one coordinate, give a folded point
$p\in\{0,\ldots,w\}$ its boundary type:
\[
 \begin{cases}
  (\mathrm{left},p),&p<r,\\
  (\mathrm{right},w-p),&w-p<r,\\
  \mathrm{interior},&r\le p\le w-r.
 \end{cases}
\]
The cases are disjoint. For an offset $s$ with $|s|\le r$,
the folded displacement from $p$ is respectively
\begin{equation}\label{eq:five-period-displacement}
 |p+s|-p,\qquad d-|d-s|\ (d=w-p),\qquad s.
\end{equation}
These identities follow directly from the two affine branches of
$\pi_w$. They also describe links in a local box when both endpoints
are in the box. Thus, after moving the anchor to cancel the folded
base point, the local field near a site depends only on its four
boundary types. The displacement formulas do not depend on $w$.
An all-interior type gives the original, unfolded field near the
chosen anchor.

For a left or right type of depth $d\in\{0,\ldots,r-1\}$,
the number of cyclic points of that type is $1$ if $d=0$ and
$2$ otherwise. The interior count is
\[
                              2(w-2r+1).
\]
Indeed its folded coordinates are $r,\ldots,w-r$, all strictly
between the two endpoints, and each has two preimages. In four
dimensions the number of sites of a prescribed type vector is the
product of its four one-dimensional counts.

Let $\mathcal T_r$ be the finite set of four-coordinate type vectors,
and let $\tau_*$ be the all-interior type. For any function
$H:\mathcal T_r\to\mathbb R$, the preceding count proves
\begin{equation}\label{eq:five-period-counting}
 \sum_{i=0}^4 c_i\sum_{q\in(\mathbb Z/2(w+i)\mathbb Z)^4}
                      H(\operatorname{type}_{w+i,r}(q))=H(\tau_*),
 \qquad c_i=\frac{(-1)^i\binom4i}{2^4\cdot4!}.
\end{equation}
To verify this exactly, fix a type vector with $d$ interior
coordinates. Its multiplicity at side $2(w+i)$ is a polynomial
in $i$ of degree $d$. The fourth difference annihilates that
polynomial when $d<4$. For the all-interior type the polynomial
is $2^4(w+i-2r+1)^4$, whose fourth difference is $2^4\cdot4!$.
Division by that factor gives one. Summing over the finitely many
type vectors proves \eqref{eq:five-period-counting}. In particular,
no boundary term has been discarded before cancellation.

\subsection{Compatible local densities}
Let $\Lambda$ be the site set of the large torus. Write
$f_\Lambda(o,U)$ for a density on the large torus and
$f_i(q,V)$, $0\le i\le4$, for densities on the five auxiliary tori
of sides $2(w+i)$. The compatibility conditions needed here are
explicit: the auxiliary densities transform covariantly under signed
lattice symmetries, the values of any two of them on globally
healthy fields agree whenever their centered radius-$r$ link
patterns agree, and the same agreement holds between each of them
and $f_\Lambda$ on globally healthy fields. Finally
$f_\Lambda(o,\cdot)$ is local on the radius-$r$ neighbourhood of
$o$.
The large period is at least $4r+16$. These are compatibility of
local functions at fixed coupling and fixed path; no equality of
partition functions on different periods is required.

Set $\widetilde f_i(q,V)=f_i(q,V)-f_i(q,\mathbf1)$ and define
$\widetilde f_\Lambda$ in the same way. The compatibility properties
persist after this centering: the identity fields are healthy and
have identical local link patterns. Define
\[
 g(o,U)=\sum_{i=0}^4c_i\sum_{q\in\Lambda_i}
                         \widetilde f_i(q,\Phi_{w+i,o}U),
                   \qquad \Lambda_i=(\mathbb Z/2(w+i)\mathbb Z)^4.
\]
Choose $z\ge w+4$. The function $g(o,\cdot)$ reads only the
radius-$z$ neighbourhood of $o$. For each auxiliary site $q$ set
$o_{i,q}=o-\pi_{w+i}(q)$, coordinatewise, and put
\begin{equation}\label{eq:five-period-shifted-localization}
 h(o,U)=\sum_{i=0}^4c_i\sum_{q\in\Lambda_i}
  \mathbf1_{\{U\text{ healthy on }N_z(o_{i,q})\}}
                  \widetilde f_i(q,\Phi_{w+i,o_{i,q}}U).
\end{equation}
The shift ensures that the physical point corresponding to $q$
is $o$. Here $N_z(a)$ denotes the radius-$z$ lattice neighbourhood
of $a$.

If $U$ is healthy on $N_{2z}(o)$, every indicator in
\eqref{eq:five-period-shifted-localization} equals one:
$\operatorname{dist}(o,o_{i,q})\le w+i\le z$ and
$N_z(o_{i,q})\subseteq N_{2z}(o)$. The folded auxiliary fields
are then globally healthy. A folded plaquette is either an
original plaquette with an orientation change, a conjugate of
such a plaquette, or a degenerate loop with identity holonomy.
The plaquette chord is invariant under conjugation and inversion;
thus its bound follows from the healthy large-field bound on the
read set.

Reflection covariance first moves $q$ to its canonical folded
representative. The displacement formulas
\eqref{eq:five-period-displacement}, followed by compatibility,
show that the summand depends only on its type vector. For the
all-interior vector it equals $\widetilde f_\Lambda(o,U)$. One can
make the last comparison using a globally healthy extension on the original large torus.
In each cyclic coordinate of length $c\ge4r$, take the track
that increases from $0$ to $2r$ in $2r$ steps, decreases back to
$0$ in the next $2r$ steps, and remains at $0$ for the remaining
$c-4r$ steps. Place its image at $o-r$ and center its argument at
the midpoint of the increasing segment. The resulting field agrees
with $U$ on the radius-$r$ box about $o$ and reads only that box
everywhere. Each plaquette is again an original plaquette with
orientation changed or an identity loop, so the extension is
globally healthy. Compatibility applies to this field on the
original period; locality then identifies its $f_\Lambda$ value with
$f_\Lambda(o,U)$. Applying \eqref{eq:five-period-counting} gives
\begin{equation}\label{eq:five-period-healthy-identity}
                h(o,U)=\widetilde f_\Lambda(o,U)
                          \quad\text{if $U$ is healthy on }N_{2z}(o).
\end{equation}

The compatibility in the selected construction is established by
passing two given periods to their product period. Replicate each
field to that period; the short local density is preserved by
replication. Translate both base points to zero. Their local
patterns still agree there, so locality of the common-period
density gives equality. The selected auxiliary representations
use the same path, coupling, profile and finite cutoff. It is
these equalities that permit the replication comparison. The
required replication and local-generation identities are part of
the short-density construction; this argument does not assert
compatibility for independently chosen histories. For the selected
densities they are proved in
Section~\ref{sec:selected-nonlinear-offering}, where the five
auxiliary densities are formed with the physical path, incoming
budget, coupling, scalar correction and finite label cutoff. When the
present section is applied, $f_\Lambda$ and $f_0,\ldots,f_4$ are
those densities; the four compatibility conditions above are then
conclusions of that section and not additional assumptions.

\subsection{The exact additive decomposition}
Let $F_a(U)$ be the ordinary logarithmic entries retained in a
finite short-support panel. Their allocation coefficients
$\omega_{oa}$ are nonnegative, sum to one in $o$ for each $a$,
and satisfy $S_a\subseteq N_r(o)$ whenever $\omega_{oa}\ne0$. Put
\[
 E(o,U)=\sum_a\omega_{oa}
           \mathbf1_{\{U\text{ healthy on }S_a\}}F_a(U).
\]
Its sum over $o$ is exactly the masked contribution of the panel.
The unmasked short density is
$f_\Lambda(o,U)=\sum_a\omega_{oa}F_a(U)$.
Translation covariance makes $f_\Lambda(o,\mathbf1)$ independent of
$o$.
Define the discrepancy
\[
                     d(o,U)=E(o,U)-f_\Lambda(o,\mathbf1)-h(o,U).
\]
This definition, followed by the change of variable
$o\mapsto o-\pi_{w+i}(q)$ in each finite site sum, proves
\begin{equation}\label{eq:five-period-exact-decomposition}
 \sum_a\mathbf1_{\{U\text{ healthy on }S_a\}}F_a(U)
 =|\Lambda|f_\Lambda(0,\mathbf1)
   +\sum_o\mathbf1_{\{U\text{ healthy on }N_z(o)\}}g(o,U)
   +\sum_o d(o,U).
\end{equation}
Indeed this change of variable identifies $\sum_oh(o,U)$ with
the middle sum. It does not require that $U$ itself be healthy.
If $U$ is healthy on $N_{2z}(o)$, the allocation supports near
$o$ are healthy, so $E(o,U)=f_\Lambda(o,U)$. Equation
\eqref{eq:five-period-healthy-identity} then gives $d(o,U)=0$.
Thus the discrepancy is supported on an unhealthy enlarged
neighbourhood. The regular output retains $g$, the constant
contribution is extracted, and the forced correction uses
$e^{-d}-1$. The negative sign follows by solving
\eqref{eq:five-period-exact-decomposition} for the new logarithmic
sum in the density convention of the nonlinear step.

\subsection{The derivative cost of folding}
Let the derivative order be $j$, the spatial weight power be $m$,
and the small-direction threshold be $\theta>0$. Suppose a
function $F$ on an auxiliary torus is $C^j$ in local exponential
charts at every globally healthy field, and has there all iterated
left derivatives through order $j$ bounded by $q$ when each direction
has norm at most $\theta$. On a healthy radius-$z$ neighbourhood, the
corresponding
bound for $F\circ\Phi_{w+i,o}$, with directions allowed to grow
as $\theta(1+\operatorname{dist}(o,x))^m$, is
\begin{equation}\label{eq:five-period-derivative-cost}
                          2^j(1+z)^{mj}q.
\end{equation}
The factor two per derivative is needed because a reversed link
transforms a left variation into a right variation. We justify
that point rather than identifying the two derivatives.

On a forward auxiliary link the induced variation is a left
variation by the corresponding original direction; on a reversed
link it is a right variation by the negative original direction.
The two types have disjoint link coordinates, fixed by the folding
map. Consequently left and right variations on different types
commute. For a list of $k$ derivatives, expand each variation into
its left and right part. There are at most $2^k$ terms. In each
term, right derivatives can be moved past all left derivatives,
and at the fixed field converted to left derivatives by adjoint
conjugation. Their norms are unchanged because the group metric
is adjoint invariant. Explicitly, start with the empty list. For a leading mixed
direction $(v_L,v_R)$ and each converted list $s$ for the remaining
directions, the two new lists are $v_L::s$ and
$s\mathbin{+\!+}[\operatorname{Ad}_{V}v_R]$, with the adjoint
chosen to turn the right action at the fixed field $V$ into a left
action. Thus the right directions occur in reverse order. The
disjoint link supports ensure that this recursion introduces no
derivatives of a varying conjugating field. Each resulting direction has norm at most
$X\theta$, where $X=(1+z)^m$. Multilinearity and rescaling by
$X^{-1}$ bound the term by $X^kq$. Summation and $k\le j$ prove
\eqref{eq:five-period-derivative-cost}. The smooth chart property
under composition follows since a link read, its inverse and a
constant identity link are smooth maps.

It follows that the ordinary localized density has derivative
price
\begin{equation}\label{eq:five-period-local-price}
       \mathcal O_z(q_0,\ldots,q_4)
                =\sum_{i=0}^4|c_i|\,2^j(1+z)^{mj}q_i,
\end{equation}
when $q_i$ bounds the centered total auxiliary density
$\sum_q\widetilde f_i(q,\cdot)$. The support is $N_z(o)$.
For the ordinary derivative budget, each entry has an assigned
anchor; its budget at a root $o$ sums only entries assigned to $o$,
weighted by $e^{e c_B(S)}$. Thus if $c_B(N_z(o))\le C_z$, the
entry $g(o)$ costs $\mathcal O_z(q)e^{eC_z}$ in that anchored
budget. A sum over every support containing a fixed site would
instead cost the additional factor $(2z+1)^4$.
A fixed signed-orientation average and its local allocation
multiply this price by the already purchased orientation cost.
All these factors depend only on the fixed local geometry, and
have upper exponent zero in the scale parameter.
